% GENERATED FILE. Edit canonical lessons, then run npm run generate:books.
% canonical-source-hash: fnv1a64:a0f3b9c8dd7043cc
% canonical-lessons: GE-C460-wetterbericht, GE-C460-wetterlage, GE-C460-wolkenlos, GE-C460-bedeckt, GE-C460-umschlagen

\chapter{Der Wetterbericht}
\label{ch:ge-der-wetterbericht}
\hlchaptermodality{460}

\begin{chapteropening}
\textbf{By the end of this chapter:} \emph{I can follow a weather report: the weather report, the weather situation, a cloudless sky, an overcast sky and a sudden change.}

Complete the last lesson of chapter 460: I can follow a weather report: the weather report, the weather situation, a cloudless sky, an overcast sky and a sudden change.
\end{chapteropening}

\section[der Wetterbericht]{der Wetterbericht — the weather report, the forecast}

\label{lesson:GE-C460-wetterbericht}



\begin{quote}
Before the new one: say the German for to graze, then the German for to milk.
\end{quote}

\subsection*{You'll want to know: der Wetterbericht}
\textbf{der Wetterbericht} — \textquotedblleft{}the weather report, the forecast\textquotedblright{}.

\begin{quote}
Laut Wetterbericht regnet es morgen den ganzen Tag.
\end{quote}

\begin{quote}
\emph{According to the weather report it will rain all day tomorrow.}
\end{quote}

Plural: \emph{die Wetterberichte}.

\begin{cousinweb}
\emph{Wetterbericht} is built from \emph{Wetter} 'weather', an inherited word and a cousin of English \emph{weather} (EtymDB), and \emph{Bericht} 'report', from Middle High German \emph{beriht} (EtymDB).
\end{cousinweb}

\begin{grammarlens}[title={What you've built}]
One more word for telling, explaining and describing at B1.
\end{grammarlens}

\subsection*{Your turn}
\begin{itemize}
\raggedright
  \item \emph{Say it:} \emph{der Wetterbericht}
  \item \emph{Say it:} \emph{der Wetterbericht}, once more
  \item \emph{Say it:} a sentence of your own with \emph{der Wetterbericht}
  \item \emph{Recall:} say the German for to graze, then the German for to milk, then say \emph{der Wetterbericht} again
\end{itemize}

\begin{tcolorbox}[breakable,colback=teal!4,colframe=teal!35!black,title={Before you move on}]
What does der Wetterbericht mean? (\textquotedblleft{}The weather report, the forecast\textquotedblright{}.) Say it once more.
\end{tcolorbox}
\section[die Wetterlage]{die Wetterlage — the weather situation, the general weather}

\label{lesson:GE-C460-wetterlage}



\begin{quote}
Before the new one: say the German for to milk, then the German for the weather report.
\end{quote}

\subsection*{You'll want to know: die Wetterlage}
\textbf{die Wetterlage} — \textquotedblleft{}the weather situation, the general weather\textquotedblright{}.

\begin{quote}
Die Wetterlage bleibt bis Sonntag stabil.
\end{quote}

\begin{quote}
\emph{The weather will stay settled until Sunday.}
\end{quote}

Plural: \emph{die Wetterlagen}.

\begin{cousinweb}
\emph{Wetterlage} is built from \emph{Wetter} 'weather' and \emph{Lage} 'position, situation', from \emph{liegen} 'to lie' (EtymDB lists \emph{Wetterlage}).
\end{cousinweb}

\begin{grammarlens}[title={What you've built}]
One more word for telling, explaining and describing at B1.
\end{grammarlens}

\subsection*{Your turn}
\begin{itemize}
\raggedright
  \item \emph{Say it:} \emph{die Wetterlage}
  \item \emph{Say it:} \emph{die Wetterlage}, once more
  \item \emph{Say it:} a sentence of your own with \emph{die Wetterlage}
  \item \emph{Recall:} say the German for to milk, then the German for the weather report, then say \emph{die Wetterlage} again
\end{itemize}

\begin{tcolorbox}[breakable,colback=teal!4,colframe=teal!35!black,title={Before you move on}]
What does die Wetterlage mean? (\textquotedblleft{}The weather situation, the general weather\textquotedblright{}.) Say it once more.
\end{tcolorbox}
\section[wolkenlos]{wolkenlos — cloudless}

\label{lesson:GE-C460-wolkenlos}



\begin{quote}
Before the new one: say the German for the weather report, then the German for the weather situation.
\end{quote}

\subsection*{You'll want to know: wolkenlos}
\textbf{wolkenlos} — \textquotedblleft{}cloudless\textquotedblright{}.

\begin{quote}
Am Nachmittag ist der Himmel wolkenlos.
\end{quote}

\begin{quote}
\emph{In the afternoon the sky is cloudless.}
\end{quote}

\emph{-los} means without: \emph{arbeitslos}, \emph{kostenlos}.

\begin{cousinweb}
\emph{wolkenlos} is built from \emph{Wolke} 'cloud', inherited from Old High German \emph{wolkan} (EtymDB), and \emph{-los} 'without', from Proto-Germanic \emph{lausaz} 'loose, free'.
\end{cousinweb}

\begin{grammarlens}[title={What you've built}]
One more word for telling, explaining and describing at B1.
\end{grammarlens}

\subsection*{Your turn}
\begin{itemize}
\raggedright
  \item \emph{Say it:} \emph{wolkenlos}
  \item \emph{Say it:} \emph{wolkenlos}, once more
  \item \emph{Say it:} a sentence of your own with \emph{wolkenlos}
  \item \emph{Recall:} say the German for the weather report, then the German for the weather situation, then say \emph{wolkenlos} again
\end{itemize}

\begin{tcolorbox}[breakable,colback=teal!4,colframe=teal!35!black,title={Before you move on}]
What does wolkenlos mean? (\textquotedblleft{}Cloudless\textquotedblright{}.) Say it once more.
\end{tcolorbox}
\section[bedeckt]{bedeckt — overcast; covered}

\label{lesson:GE-C460-bedeckt}



\begin{quote}
Before the new one: say the German for the weather situation, then the German for cloudless.
\end{quote}

\subsection*{You'll want to know: bedeckt}
\textbf{bedeckt} — \textquotedblleft{}overcast; covered\textquotedblright{}.

\begin{quote}
Morgens ist der Himmel bedeckt, später kommt die Sonne.
\end{quote}

\begin{quote}
\emph{In the morning the sky is overcast; later the sun comes out.}
\end{quote}

The past participle of \emph{bedecken} 'to cover', used as an adjective.

\begin{cousinweb}
\emph{bedeckt} is the participle of \emph{bedecken} 'to cover', from \emph{decken}, inherited from Old High German \emph{decchen} and Proto-Germanic \emph{þakjaną} (EtymDB), the root of \emph{Dach} 'roof'.
\end{cousinweb}

\begin{grammarlens}[title={What you've built}]
One more word for telling, explaining and describing at B1.
\end{grammarlens}

\subsection*{Your turn}
\begin{itemize}
\raggedright
  \item \emph{Say it:} \emph{bedeckt}
  \item \emph{Say it:} \emph{bedeckt}, once more
  \item \emph{Say it:} a sentence of your own with \emph{bedeckt}
  \item \emph{Recall:} say the German for the weather situation, then the German for cloudless, then say \emph{bedeckt} again
\end{itemize}

\begin{tcolorbox}[breakable,colback=teal!4,colframe=teal!35!black,title={Before you move on}]
What does bedeckt mean? (\textquotedblleft{}Overcast; covered\textquotedblright{}.) Say it once more.
\end{tcolorbox}
\section[umschlagen]{umschlagen — to change suddenly (of the weather); to turn over}

\label{lesson:GE-C460-umschlagen}



\begin{quote}
Before the new one: say the German for cloudless, then the German for overcast.
\end{quote}

\subsection*{You'll want to know: umschlagen}
\textbf{umschlagen} — \textquotedblleft{}to change suddenly (of the weather); to turn over\textquotedblright{}.

\begin{quote}
In den Bergen schlägt das Wetter schnell um.
\end{quote}

\begin{quote}
\emph{In the mountains the weather changes quickly.}
\end{quote}

Separable; Perfekt with \emph{sein} for the weather: \emph{das Wetter ist umgeschlagen}. \emph{Der Umschlag} — the envelope.

\begin{cousinweb}
\emph{umschlagen} is built from \emph{um} 'over, round' and \emph{schlagen} 'to strike, to turn', inherited from Proto-Germanic \emph{slahaną} (EtymDB).
\end{cousinweb}

\begin{grammarlens}[title={What you've built}]
One more word for telling, explaining and describing at B1.
\end{grammarlens}

\subsection*{Your turn}
\begin{itemize}
\raggedright
  \item \emph{Say it:} \emph{umschlagen}
  \item \emph{Say it:} \emph{umschlagen}, once more
  \item \emph{Say it:} a sentence of your own with \emph{umschlagen}
  \item \emph{Recall:} say the German for cloudless, then the German for overcast, then say \emph{umschlagen} again
\end{itemize}

\begin{tcolorbox}[breakable,colback=teal!4,colframe=teal!35!black,title={Before you move on}]
What does umschlagen mean? (\textquotedblleft{}To change suddenly (of the weather); to turn over\textquotedblright{}.) Say it once more.
\end{tcolorbox}
